\chapter{Statistical physics}

It is not surprising that the modern landscape, insofar of analysis and deep learning architectures, inspired a particular cohort of researchers with ideas on using \textbf{statistical physics} to infer on a learning theoretic system as for it to be a \textit{statistical physics-induced complex system analysis}, i.e. a complex multi-degree of freedom system that statistical physics can take a hold on. The other direction is to apply artificial neural network to simulate a complex interacting systems, and thus verifying a version of mathematical descriptions on it, such as the \textbf{duality} of statistical physical descriptions.  The list of principle illustrative works on this is not lacklustre in the least, such as
\begin{enumerate}[topsep=1pt,itemsep=1pt]
    \item \cite{ferrari2024machinelearningoptimizationbasedapproaches}: Machine learning and optimization-based approaches to duality in statistical physics, i.e. where two mathematical descriptions can describe the same system. The paper used artificial neural network, which benefits from arbitrary component encoding, thus allow for representation on arbitrariness. Similar analysis is used of ML-assisted system in \cite{hou2023machinelearningrenormalizationgroup}, for the renormalization group. 
    \item \cite{barbier2025statisticalphysicsdeeplearning}: Works on the statistical physics of deep learning, and sought to analyse deep neural network similar to statistical system, and thus derive an optimal learning scheme on such. Works similar to this can be of comparison, toward \textbf{statistical physics interpretation}, such as \cite{pei2025statisticalphysicsartificialneural,Faulkner_2024}, or specialized case, most often as \textbf{Graph Neural Network} (GNN) like in the case of \cite{Duranthon_2025}. An interesting phenomenon or observation on deep learning system, and marginally classical learning systems, is \textbf{phase transition}, which is explored in \cite{ziyin2022exactphasetransitionsdeep} for example, and \cite{PhysRevResearch.5.043243} for a recent concept of first-order and second-order understanding of deep neural network interpretation for phase transitions. 
\end{enumerate}

Utilization of such statistical physics structures onto deep learning includes \cite{wainwright_jordan_2008} of the Ising or Markov Random Field models, Sherrington-Kirkpatrick (SK) spin glass in \cite{sherrington_kirkpatrick_1975} for nonconvex loss landscape in ML (as per direct interpretation language); similar purpose is presented by using \textbf{p-spin or spherical p-spin models} (\cite{crisanti_sommers_1992}). Specifically on neural network and deep learning, \cite{hopfield_1982,ackley_hinton_sejnowski_1985,hinton_2002_cd} and \cite{gardner_1988} utilized many statistical physics analogues onto neural networks, though in Gardner case it is a \textbf{perceptron optimal-capacity model}. In modern time, it is the \textbf{cavity model} on high-dimensional understanding, as in \cite{donoho_maleki_montanari_2009}, \cite{Jacot:2018:NTK} Neural Tangent Kernel (NTK) with statistical inspirations, Langevin mechanics of stochastic dynamics analogue as in \cite{mandt_2017}. Energy-based models, as in \cite{ho2020ddpm,song2019score} and entropy-SGD of \cite{chaudhari_entropy_sgd_2016}, mean-field analysis (\cite{mei_montanari_2018}), and a rather modern \textbf{free-energy interpretation} of modern ML, in \cite{kingma_welling_2013}. It is fair to say that statistical physics abstraction itself gives plenty applications and inspirations on the current modern ML landscape. 

However, it is our subjective opinion that is to validate, that statistical physics application, and most of its concurrent analysis are inherently messy, disorganized, and of intrusive mapping from statistical physics onto deep learning or classical systems. We notice a central pattern that typically, statistical analogues are applied onto the machine learning or dynamic model system, without cares and consideration for a direct \textit{formal and conceptual mapping} between the two theories. For example, it is not sure how well all different axiomatic presumption of statistical physics, statistical field theory works on machine learning system, or simple as old classical models with human-aided designs in the 1970s-1980s, or symbolic branches of machine learnable methods. Concepts like temperatures and entropy are typically presented as fact, or rather a list of intrinsic properties gave rise to a `natural measure', yet is not so applicable to a wider degree of the deep learning landscape. Interpretation exploded during the rise of deep learning models, as \cite{goodfellow2016deep} also introduced some in his book, since neural network produces a corridor for \textit{isolated component structures}, treated similar to statistical physics without structural difference between the two. As such, there exists the decoupling, that leads to Hinton's both `failure' and `success' of the Boltzmann neural network (\cite{ackley_hinton_sejnowski_1985,hinton2012practical}), being a failure since it requires to express neural network in a niche, toy-liked, explicitly redesigned, narrow structure, and closed end-to-end loop for equilibrium evolution in-system. Such is often considered to be utilizing a representation scheme, this time, neural network, with fixed descriptions on encoding choice, rather than neural network in general. 

Thus, we lay our goal to identify this gap, clarify concepts and deliver a closest possible working the problem, of applying statistical physics scheme onto the ontologically different space of objects. Within reasons, we also utilize such onto the set of \textbf{double descent problem}, in which way interpretations make sense and historical patterns or theories, and some experiment to test the hypothesis. 

\section{Note for team writing member}

Please skip this page, right of the \texttt{clearpage} command right behind this. It is for organization during paper writing, and will be removed afterward. 

\subsection{Main structure}

The following outlines the organizational structure in proposal specifically for the statistical physics paper. There are two parts, few claims that we posit in such discussion. Those are the main central point, but can also be mutated when further evolutions are possible. 

\begin{enumerate}[topsep=1pt,itemsep=1pt]
    \item We posit that statistical physics is a quite formal and general method, yet unrefined, and often is applied incorrectly in the framework of perspective of machine learning. There are certain situations (or rather frequently) that statistical physics do describe a deeper interpretation of a dynamical system, but on the other hand, forces on machine learning the \textbf{physical perspective} and description, and thus enables false claims, restricted claims, setup that is impossible, of assumptions that never satisfy the theory itself, to force everything into a physical system instead of machine learning, then claiming the reverse is the truth. 
    \item We claim that to understand statistical physics interpretation to machine learning or system, we need more than apply superficially, that is, to simply do the one-one mapping of this framework's concept to another. 
    \item We will highlight such mismatch in the analysis, both in details of interpretation, and both formal way of mathematical expression, proofs, or so on, concretely. 
    \item We will propose a way to look at such, particularly through the problem of double descent, where other theories or statistical physics interpretation fails, and the analogous similarity to other previous theories and attempts historically. 
    \item We then bring up forth a new theory, or rather, a new abstract framework to explain it. Couple with that are some toy experiments, to explain of the system we want to highlight. 
\end{enumerate}

Thus, the following sectioning is adopted as the formal form for the document. This includes sections on the introduction of statistical theory, the argument and concurrent (present theory) analysis, then reformulation, then transition to a new theory with experiments. 

\begin{enumerate}[topsep=1pt,itemsep=1pt]
    \item Introduction to Statistical Physics: Simply the introduction of the discussion about the main position, as to what is statistical physics, and the morphism to machine learning, or the learning theory. 
    \begin{enumerate}[topsep=1pt,itemsep=1pt]
        \item Reasons of statistical physics. This is to dissect \textit{what}, and why statistical physics is even applied, and what does it mean or how to specifically connect to machine learning. This includes the analysis of different \textit{abstraction} level that must be made to apply certain parameters or so on of ML, to fit statistical physics (for example, what is required to \textit{simplify} or \textit{add details} into it such that the notion of temperature is borderline acceptable in ML algorithms?).
        \item Axioms of statistical physics. What are the axioms and formulation that statistical physics use, from assuming properties of multi-particle interpretation, the dynamic allowed, what kind of dynamics internally and externally, etc., such that statistical physics is expressed most generally? And what can be of the application and mismatch of between? 
        \item Assumptions of statistical physics. Explicit assumptions list, and incompatibility explanation, derivation, analysis. For example, it can be like that independent of components system i.e. the molecular behaviours, which is often not true in an integrated ML system. 
        \item Relevant concepts that it uses in formation, assumptions, already established notions. 
    \end{enumerate}
    \item Problem and settings. Here, it is simple as to \textit{identify the errors}, speaking out loud of the problems, the problems of people solving the problems, the fallacy that make everything fails, what can be retained, what actually works, its limitation and so on. 
    \item Identification: Strictly identify the problem solutions and identification on "what exactly the problem can be categorized", proposing hypothesis, solutions, there of. 
    \item Reformulation: Reformulate it using a new theory. Chooses two mainline approach. One on the \textbf{predictive patterns}, and one on the \textbf{structural formation}, i.e. not just predicting \textit{the peak or rise} in double descent, but to identify the structural reason why that happens, and so on. 
\end{enumerate}

All of this works in conjunction with double descent as a medium of illustration. Specifically, the \textit{problem and setting} part would have to utilize in-context illustration, which then would be the double descent experimental cases. So is \textit{identification} and reformulation. The \textbf{reformulation} part is however much more general, but use double descent as a toy problem for test bed. 

\subsection{Notices}

There are a few points that must be addressed. This will help and will be included in a much more sophisticated writing style prefix than before, ensuring that details will be added now and not later where necessary. 

\begin{enumerate}[topsep=1pt,itemsep=1pt]
    \item \textbf{Results are written as if they are already settled theory}. The entire thing contains intermediate reasoning, failed paths, test cases, and scope limits are erased (actually provided later but still, it is \textit{not finalized but pushed as if is}). You need \textbf{expendable structure} instead of \textbf{finished fragmented endpoints}. 
    \item The current working of the statistical physics treatment below is \textit{suboptimal}. You jump to the conclusion too fast, and whenever submissions are in, it is in the state of already finalized semi-manuscript with conclusion. Elaborate on details, testing out cases, and so on, and make incremental gains as cumulative on a main document, not simply `separated documents' that cannot be used in tangent with other team members. 
    \item Hiding too many stuffs. For example, even though it is said that there is isomorphism, you \textbf{must} clarify what is isomorphism here, why it is here, what is it, what isomorphic to what. All concepts are the same. Phase transition? If you list it out or seriously consider it part of the theory, include it or explain in great details. 
    \item Even though you stated, for example, the \textbf{ergodic theorem} or axioms or principle, you must \textit{explicitly explain them}, with \textbf{full setting} of what object to consider, what is the central lining of the axioms, what it asserts, why is it there, what situation makes it to not \textit{work} any more, and vice versa, add many more of the full details as to then justify the content itself. 
    \item Identify \textbf{actual notational system} and how they mean, as ontology and notations must be clarified of for one symbol contains many contexts and so on, especially with a long manuscript. Even Langevin dynamics will have to be clarified as to \textbf{why it is defined like that}, and thus why machine learning people use it in the theory. Exploit many topics, for example, Hamiltonian of a system but now ask `since we are working in machine learning, why don't we test how and who the heck did the Hamiltonian of a ML system, or are there any \textit{analogous} notion in here?' so on an so forth. Equilibrium assumption is very \textit{structural}, and requires many knowledges beforehand, so \textbf{list them out carefully and elaborate}. 
\end{enumerate}

Overall, this section misses details. It misses the basic foundational sections, explanation, completeness of details in a sensitive setting that we \textbf{interpret} and \textit{unify} a fragmented connection between two fields, of which one assumes higher authority (statistical physics) than the other (learning theory). If we keep doing this, everything will fail and misrepresentation will happen, overclaiming will happen, something that seems to be important and big will turn out to be a slight cheeky representation and nothing more. And all the good ideas might just be not developed in blindness of the other mediocre but `nice-looking' and hand-wainable. 

This is a setback and must be addressed. The concept itself is highly valuable in the followed idea templates, and frankly to say genuinely good, but the elaboration and formation of such is lacking seriously in several fronts. Also, the \textbf{styling} is also too `short' and is too dense in the wrong way of dense - in the sense of compacting many technical jargons but confusing in many parts. For example, `energy landscape as rugged, multi-valleyed, dominated by metastability' contains the concept of a \textbf{rugged affine $n$-hyperspace in a $n+1$ space},  convex peak or valley and \textbf{why is it even there} on the metric of the system, metastability but does not clarify whether it is a \textbf{global metastability} from chaotic distribution or something else. Spin-glass theory is also non-trivial, but is sighted without prior blanket. This standard \textbf{must be met} without question. Another thing to finally add as to why that kind of density is dangerous is that, from that simple sentence, you have several shots onto your leg:

\begin{enumerate}[itemsep=1pt,topsep=1pt]
    \item What metric defines distance in the $n$-space? What is that space and the structures provided? Does it have inner product? Why the heck inner product is even in there? Suddenly now a configuration space is given pseudometric - why? 
    \item What notion of convexity is being used? What if it is not convex? What is the reference point (usually convexity is defined by 'probing two points' intuitively speaking)? What is alternative than convex? Why is convex there? Is the theory restricted to this convexity or that convexity and if it is non-convex in general, is it useless? Is convexity even a principle subject that defines the theory?
    \item Is ruggedness quantified (e.g. Hessian spectra, barrier scaling)? What is rugged? Under what context? Is it emergent or inherent of a system up to certain threshold?
    \item Is metastability dynamical (algorithm-dependent) or thermodynamic? What is being thermodynamic here means? Is it about a thermodynamic interpretation of a componential deep neural network structure? Is the algorithm being agnostic? Is the algorithm of different nature (physical operational algorithm or simply numerical or automata-kind one)?
    \item What is the embedding space, what is encoded, how it is encoded, what is encoded? Even if information theoretic ideas are imported in interpretation from us researchers, \textbf{what does it mean for a system to have information?} Is it intrinsic, or observers' interpretation only? 
\end{enumerate}

Too many vector of analysis and thus too many fall-line that makes any further attempt to collapse under its own weight. We are not in the stage of building ideas any more in a fragmented way. There are guidelines which I provided, and there are structural in which we agreed upon. Ideas go in ideas, but structural must stay loyal to the structure and not to avoidance of details. The manuscript is already 150 pages long. There is nothing to shy from \textit{keeping it short either} as the nature of this document. This is for anyone, so do \textbf{myself} on this section that I keep in my mind:

\begin{quote}
    You can't expect people to understand what you write when you yourself cannot even invoke it whenever asked. Don't delegate meaning to those you subconsciously know, because you also know consciously you don't. In foundational, unifying, or interpretive work, the opposite is true. Every unstated assumption \textbf{becomes a liability}, every borrowed term \textbf{must be re-earned}, and every failure mode \textbf{must be named}.
\end{quote}

Keep on with this and it would be alright. Otherwise, foundational work will find its foundation in another corner of a dusty library and nothing more. If they cannot meet this standard, they are free to leave and fragments of what they have done will still be recognized in the manuscript as proof of their involvement, no question asked. However, this is imperative. They must adhere to this and so am I, also free to be criticized when wrong. Science does not await for anyone to be soft-talked into obscurity. When I asked of you this project being foundational, I said to beware of the contract you must sign before you commit. Whether you now feel it too harsh of your own or not, the judgement is yours. But if you still work, this will be the judgement I give before you improve on this, for the betterment of both our work and ourselves.

\subsection{In detail systematic projection}

Now then, for the main portions, we should go with the following list, though it can be modified any given time. If you find something lacking, please insert it in this frame on details more than the above conceptual schematic, providing it is a good addition to the fold. In general, there is
\begin{enumerate}[topsep=1pt,itemsep=1pt]
    \item Introduction to statistical physics. Specifically, this part will grapple about the \textit{foundational criteria of statistical physics}, i.e. what is the system of laws, how objects are presented. This includes the following topic to include in, not per specific order but relative conceptual order:
    \begin{itemize}[topsep=1pt,itemsep=1pt]
        \item What is statistical physics? Why is it necessary?
        \item What is the foundational framework that precedes it? I.e., for example, statistical physics of course goes before quantum but still applicable in quantum system with modification, so it effectively \textbf{switch register} in between. So, let's say we go to the foundation on how Newtonian foundation, theoretical mechanics and the needs of statistical method influences such foundation and thus the assumptions that follow. 
        \item The tools and perspective that statistical physics uses for such problem and so on. Especially on, for example, the \textbf{simulated scale problem} that statistical physics seems to work around, because of the requirement to \textit{estimate large and multi-body} systems.  
        \item What is the tradeoff in forming that system, what is assumed during the process of `statisticalized' the system and a new second-layer ordering of axiomatic set, for example? We need to systematically identify this, by proposing your own questions, interpretations, aspects. 
        \item Now, after we are done, move on to the typical and standard analysis, layout and rigorous analysis of statistical mechanics in actions, the principles, the theorems, the system analysis and observations in the form of physical theory that statistical theory has made. Going for this into the heart of the theory itself, especially those that would often be used in ML, like replica methods, and why it is there, is a crucial problem. 
        \item Following the previous step, ask - why you use \textit{replica method} and so on, but not the others' analysis? This is because theory in itself compounds from analysis and aspects, or descriptions and encoding in-field, but this is where you deliberately \textbf{map the theory to another place}, which means that this particular relationship is rather `toxic' - the other might not have what you previously identify as \textit{natural evolution in solution of methods}. For example, even the action of formalizing Newtonian mechanics to Lagrangian requires the presumption and properties by uniformity of system description - meta-theoretical duality, so to speak.  
    \end{itemize}
    \item The connection to learning theory, and the modelling schema in general. This includes the following problem sections:
    \begin{itemize}
        \item Current methods and categorization of them (with references as always) from statistical physics and was used in ML. Similarly, analogous concepts from statistical physics to ML, and the reverse relationship if any in the way back. 
        \item Present results and historical researches, caveats in great details, analysis of the results themselves, author's own observations, the scope limitation, the setting self-constraints, and so on. 
        \item Analysis of each fundamental case - what is used, what is the \textbf{object of concept} in statistical physics, and \textbf{object under transformation} of understanding in learning theory or computational model used, i.e. what is the core objects and so on that are mapped to \textit{what object over there} in ML,TML, learning theory and so on. 
        \item Amplification on detail research of drawbacks, i.e. this is the section to limit the scope and range of applicability of those above concepts and researches, nitpick on them, and so on. The goal is to reframe their reach here, and to provide windows of opportunity resulted from such act. 
        \item Replication of old experiments and applications, and explore their specific edge cases. 
    \end{itemize}
    \item Double descent on learning models and the statistical compounding explanation. This section is rather straightforward - find out what is going on with double descent and \textbf{the statistical physical explanation already presented}. This is crucial as to connect back to earlier sections on \textit{why it works} and when it does not, and to analyse the specific double-descent specific categories here. That said, focus on the conclusion or methods used on the notion of \textbf{model complexity}, \textbf{model structure}, \textbf{error landscape and model evaluation}, then \textbf{algorithmic performance and operational awareness} (i.e. the difference between deterministic and probabilistic plus operational systems are that they are volatile enough for us to justify actual operational features), \textbf{algorithm complexity} both in the time-complexity and structural complexity, agnosticity of algorithm and operational processes of models, hidden assumptions and reduction of scalability, for example, in certain texts, they specifically \textit{constraints} the appearance of double descent to a very niche portion of the domain - if they generalize it further than it should be - that is our cue. Additionally, if able, the structural changes and pathological pathways explored by those models. 
    \item The novel development. I will leave this to you and the collaborators else in your own judgement. However, as said, there are two types of theoretical development in question - the \textit{predictive framework} and the \textit{structural framework}. The first one succeed in \textit{the prediction is inherent of this domain and system}, while failing at to point out specifically why and which components break, while the other one \textit{can explain well the structural reasons that lead to this particular phenomenon}, but phenomenologically performs worse do to the nature of structural conformation.
\end{enumerate}

The general inquiry of this is to develop them in rigorous detail, for the repeated usage of this particular word as-is. We do not require being super-formal and mathematical, but the right way of being rigorous - as to hide nothing, unpack everything, precise language or can be often generating as the topic itself is largely uncharted domain, and mathematically \textit{effective} - we don't want performative mathematical overformalism. The writing style fixture has been included in earlier notes, so we should not bring it here. However, if to add, then the mainline argument is to "don't overclaim where shouldn't, overcommit where trivial, under-communicate and describe where obscure, attack persons where should be targeted toward their works". 

\section{Approach A - Isomorphism}

Isomorphism here, for the sake of clarify means the \textbf{structure preservation mapping} between two systems of abstract objects. The categorical error here is to assume isomorphism between two fields, and mapping them onto the general line of basis, and is a very famous \textit{failure point} in the entire thesis of this section. 

\subsection{From Analogy to Isomorphism}

IN here, we critically examine the intersection of statistical physics and machine learning. We posit that while statistical physics provides a rigorous framework for analyzing complex systems with many degrees of freedom, its application in machine learning has often been hindered by superficial analogies rather than mathematical isomorphisms. To truly leverage this framework, one must move beyond metaphorical comparisons and establish a dictionary of strict correspondence between physical quantities and learning dynamics.

\subsubsection{Reasons for Statistical Physics in Machine Learning}

The primary motivation for applying statistical physics to machine learning lies in the ``curse of dimensionality'' and the emergent phenomena observed in deep neural networks. Just as thermodynamics describes the macroscopic behaviour of $10^{23}$ particles without tracking individual trajectories, statistical learning theory seeks to describe the generalization performance of models with billions of parameters. However, a critical mismatch often occurs in the abstraction level.

Common approaches often equate the ``learning rate'' in Stochastic Gradient Descent (SGD) to ``temperature'' in Langevin dynamics. We argue that this abstraction is insufficient and often misleading. In equilibrium physics, temperature is an isotropic scalar field. In contrast, the noise introduced by SGD is inherently anisotropic and state-dependent, scaled by the gradient covariance matrix. Therefore, a direct application of equilibrium concepts (like canonical distributions) without accounting for this structural mismatch leads to incorrect predictions about the system's stationary state.

To be more precise, let us consider the standard Langevin dynamics:
\begin{equation}
    dW = -\nabla L(W) \, dt + \sqrt{2T} \, dB_t
\end{equation}
where $T$ is the temperature and $dB_t$ is a Wiener process. The stationary distribution is the Boltzmann distribution $\rho(W) \propto \exp(-L(W)/T)$. However, SGD dynamics are fundamentally different:
\begin{equation}
    W_{t+1} = W_t - \eta \nabla L_{\mathcal{B}}(W_t)
\end{equation}
where $\mathcal{B}$ is a mini-batch. The noise covariance is:
\begin{equation}
    \Sigma(W) = \mathbb{E}_{\mathcal{B}}\left[(\nabla L_{\mathcal{B}} - \nabla L)(\nabla L_{\mathcal{B}} - \nabla L)^\top\right]
\end{equation}
This covariance is \textbf{state-dependent} and \textbf{anisotropic}, fundamentally different from isotropic thermal noise.

\subsubsection{Axioms and Formulation}

Statistical physics is built upon fundamental axioms, most notably the \textit{Ergodic Hypothesis}, which assumes that a system evolves over time to explore all accessible microstates with a probability proportional to their Boltzmann weight.

\paragraph{The Mismatch.} Deep learning dynamics, particularly under SGD, are fundamentally non-ergodic on practical timescales. The system does not explore the entire weight space but is confined to specific trajectories determined by the initialization and the landscape geometry.

\paragraph{The Isomorphism.} To resolve this, we propose a strict isomorphism:
\begin{itemize}[topsep=1pt,itemsep=1pt]
    \item \textbf{Microstate:} A specific configuration of weights $W \in \mathbb{R}^P$.
    \item \textbf{Hamiltonian (Energy):} The Loss function $L(W, \mathcal{D})$.
    \item \textbf{Quenched Disorder:} The dataset $\mathcal{D}$, which remains fixed during the training dynamics, analogous to the random interactions $J_{ij}$ in spin-glass models.
\end{itemize}

This mapping situates machine learning not in the realm of simple ideal gases, but squarely within the physics of \textit{disordered systems}, where the energy landscape is rugged, multi-valleyed, and dominated by metastability. The disorder average over datasets corresponds to the quenched average in spin-glass theory:
\begin{equation}
    \overline{\langle O \rangle} = \int d\mathcal{D} \, P(\mathcal{D}) \langle O \rangle_{\mathcal{D}}
\end{equation}
where $\langle \cdot \rangle_{\mathcal{D}}$ denotes the thermal average for fixed dataset $\mathcal{D}$.

\subsubsection{Assumptions and Incompatibilities}

Standard statistical mechanical treatments (e.g., the Replica Method) assume the system reaches thermodynamic equilibrium. This implies that the final solution is independent of the training path. However, the phenomenon of double descent---and specifically epoch-wise double descent---demonstrates that the ``solution'' is a transient state of a non-equilibrium process. 

\begin{assumption}[Equilibrium Assumption]
The system reaches a stationary distribution $\rho^*(W)$ independent of initial conditions and training trajectory.
\end{assumption}

This assumption is violated in practical deep learning for several reasons:
\begin{enumerate}[topsep=1pt,itemsep=1pt]
    \item \textbf{Finite training time:} Training is stopped before equilibration.
    \item \textbf{Non-convex landscape:} Multiple local minima prevent ergodic exploration.
    \item \textbf{Time-dependent dynamics:} Learning rate schedules break time-translation invariance.
\end{enumerate}

Forcing an equilibrium assumption onto a dynamic process like SGD is the fundamental error we aim to correct in subsequent sections.

\subsection{Limitations of Equilibrium Theories}

Here, we identify the specific failures of current theoretical frameworks when addressing the Double Descent (DD) phenomenon. While significant progress has been made, a comprehensive theory remains elusive due to a reliance on static analysis.

\subsubsection{The Failure of Classical Bias-Variance}

As established in Theorem~\ref{thm:bias-variance}, the classical U-shaped curve predicts that as model complexity increases beyond the optimal point, the variance term dominates and test error monotonically increases. Formally, for the decomposition:
\begin{equation}
    \mathbb{E}[\mathcal{M}(f,y)] = \mathcal{B}(f,y) + \mathcal{V}(f,y) + \sigma^2_{\text{noise}}
\end{equation}
classical theory predicts $\mathcal{V}(f,y) \to \infty$ as model complexity $\to \infty$. However, empirical evidence from \cite{belkin_reconciling_2019,nakkiran_deep_2019} shows a second descent in test error for over-parameterized models ($P > N$).

\subsubsection{The Limitation of RMT and Replica Theory}

Recent applications of Random Matrix Theory (RMT) and Replica Theory have successfully derived the ``peak'' of the double descent curve at the interpolation threshold ($P \approx N$). For linear regression with random features, the expected test error takes the form:
\begin{equation}
    R_{\text{test}} = \sigma^2 \frac{P}{N - P - 1} + \text{bias terms}, \quad P < N
\end{equation}
which diverges as $P \to N$. However, these successes are limited to ``static'' or ``shallow'' models. They treat the solution as a global minimum (min-norm solution) found in a convex landscape:
\begin{equation}
    W^* = \argmin_W \|W\|^2 \quad \text{s.t.} \quad XW = Y
\end{equation}

\paragraph{The Identification of Errors.} The critical error in these approaches is the neglect of \textit{feature learning dynamics}. Deep neural networks do not merely fit fixed features; they evolve their representation space. Furthermore, the existence of \textit{epoch-wise double descent}---where the test error exhibits non-monotonic behavior over time even with fixed model size---proves that DD is a dynamical phenomenon. Equilibrium theories, by definition, cannot explain transient dynamical effects.

\section{Identification: The Dynamic-Landscape Interaction}

We categorize the Double Descent problem not merely as a statistical anomaly, but as a complex interaction between \textit{multi-scale dynamics} and a \textit{critical landscape}.

\subsubsection{The Gap Between Statics and Dynamics}

There exists a ``gap'' between the static description of the loss landscape and the dynamic trajectory of the optimizer. We identify three distinct regimes that must be unified:

\begin{description}
    \item[Under-parameterized Regime ($P \ll N$):] Dominated by simple signal learning. The landscape is relatively smooth with a unique global minimum.
    
    \item[Critical Regime ($P \approx N$) --- Jamming:] At $P \approx N$, the system undergoes a ``Jamming Transition''. The landscape becomes rugged with sharp minima, corresponding to satisfying a precise number of constraints equal to the degrees of freedom. The condition number of the Hessian diverges:
    \begin{equation}
        \kappa(H) = \frac{\lambda_{\max}}{\lambda_{\min}} \to \infty \quad \text{as} \quad P \to N
    \end{equation}
    
    \item[Over-parameterized Regime ($P \gg N$):] The system becomes ``floppy'' or unjammed. There exists a continuous manifold of zero-loss solutions, and the geometry of this manifold determines generalization.
\end{description}

\subsubsection{Hypothesis: Double Descent as Phase Transition Navigation}

\begin{hypothesis}[Dynamic Origin of Double Descent]
Double Descent is the result of the SGD algorithm navigating through phase transitions in the loss landscape. Specifically:
\begin{enumerate}[topsep=1pt,itemsep=1pt]
    \item The ``peak'' at $P \approx N$ is a trapping event in the jammed phase (sharp minima).
    \item The second descent is a relaxation process driven by the specific noise structure of SGD, which allows the system to escape sharp minima and settle into flat, generalizing minima.
\end{enumerate}
\end{hypothesis}

The key insight is that SGD noise is not merely a nuisance but an \textit{implicit regularizer} that performs geometric selection among the manifold of solutions.

\subsubsection{Reformulation: Non-Equilibrium Statistical Physics Framework}

To address the identified gaps, we reformulate the learning problem using a Non-Equilibrium Statistical Physics (NESP) framework. We move from describing static probability distributions to modeling the evolution of the probability density $\rho(W,t)$.

\subsubsection{The Fokker-Planck Description}

The evolution of the probability density over weight space is governed by the Fokker-Planck equation:
\begin{equation}
    \frac{\partial \rho(W, t)}{\partial t} = \nabla \cdot \left(\rho \nabla L\right) + \nabla \cdot \left(D(W) \nabla \rho\right)
\end{equation}

Crucially, the diffusion tensor $D(W)$ is not constant (as in thermal equilibrium) but is \textbf{state-dependent}, defined by the covariance of the stochastic gradients:
\begin{equation}
    D(W) = \frac{\eta}{2B} \Sigma(W)
\end{equation}
where $\eta$ is the learning rate, $B$ is the batch size, and $\Sigma(W)$ is the gradient noise covariance.

\subsubsection{Noise-Driven Entropic Selection}

This reformulation leads to a new mechanism for generalization. The SGD noise is anisotropic and approximately proportional to the Hessian (curvature) of the loss landscape:
\begin{equation}
    \Sigma(W) \approx \frac{1}{N} \sum_{i=1}^N \nabla \ell_i \nabla \ell_i^\top \sim H(W) \quad \text{near minima}
\end{equation}

This leads to curvature-dependent diffusion:
\begin{itemize}[topsep=1pt,itemsep=1pt]
    \item \textbf{In Sharp Minima} (associated with the interpolation peak/jamming): The curvature is high, so $D(W)$ is large. This creates a strong effective ``repulsive force'' that drives the system out of these non-generalizing states.
    
    \item \textbf{In Flat Minima} (associated with the second descent): The curvature is low, so $D(W)$ is small. These states are dynamically stable under SGD.
\end{itemize}

\begin{figure}[htb]
    \centering
    \resizebox{0.5\textwidth}{!}{%
    \begin{tikzpicture}[scale=1, >=Stealth]
        % --- Axes ---
        \draw[->] (-0.5,0) -- (9,0) node[right] {Weight Space ($W$)};
        \draw[->] (0,-0.5) -- (0,5) node[above] {Loss / Energy $L(W)$};

        % --- Landscape Curve ---
        \draw[thick, blue, smooth, domain=0.5:8.5, samples=100] 
            plot (\x, {0.05*(\x-2)^2 * (\x-6.5)^2 + 2*exp(-10*(\x-2)^2) + 0.5});
            
        % --- Sharp Minima Elements ---
        \node at (2, 2.6) [circle, fill=red, inner sep=2pt] (sharp_particle) {};
        \draw[->, red, thick, decorate, decoration={snake, amplitude=.4mm, segment length=2mm}] (2, 2.6) -- (1, 3.5) node[midway, left, font=\scriptsize] {High Noise};
        \draw[->, red, thick, decorate, decoration={snake, amplitude=.4mm, segment length=2mm}] (2, 2.6) -- (3, 3.5);
        \node[align=center, font=\small] at (2, -1.2) {\textbf{Sharp Minima} \\ (Jamming Phase) \\ High Curvature};
        \draw[dashed] (2, 2.5) -- (2, 0);

        % --- Flat Minima Elements ---
        \node at (6.5, 0.6) [circle, fill=green!60!black, inner sep=2pt] (flat_particle) {};
        \draw[->, green!60!black, thin] (6.5, 0.6) -- (6.2, 0.9);
        \draw[->, green!60!black, thin] (6.5, 0.6) -- (6.8, 0.9) node[midway, right, font=\scriptsize] {Low Noise};
        \node[align=center, font=\small] at (6.5, -1.2) {\textbf{Flat Minima} \\ (Relaxation Phase) \\ Low Curvature};
        \draw[dashed] (6.5, 0.5) -- (6.5, 0);

        % --- Transition Arrow ---
        \draw[->, dashed, thick, purple] (2.2, 3) to[bend left=30] node[above, align=center, font=\small] {\textbf{NESP Transition}\\(Annealing)} (5.5, 1.5);
        
        % --- Annotations ---
        \node[align=left, font=\scriptsize, draw] at (6, 5) {
            \textbf{Mechanism:} \\
            Noise $\propto$ Curvature \\
            $D(W) \sim \nabla^2 L$
        };
    \end{tikzpicture}%
    }% end \resizebox
    \caption{{\small \textbf{NESP Mechanism of Double Descent.} At the interpolation threshold (Sharp Minima), high curvature induces strong anisotropic SGD noise (red wavy arrows), driving the system to escape. The system then relaxes to flat minima where low noise ensures stability. This process explains the second descent in test error.}}
    \label{fig:nesp_mechanism}
\end{figure}


\subsubsection{The Mechanism of Double Descent}

This framework reinterprets the Double Descent curve as a dynamical timeline:

\begin{enumerate}[topsep=1pt,itemsep=1pt]
    \item \textbf{Descent 1 (Signal Learning):} Fast learning of dominant signal components (low-frequency modes). The model captures the primary structure in the data.
    
    \item \textbf{Ascent (Peak/Jamming):} The system attempts to fit high-frequency noise but encounters the ``Jamming transition'' at the interpolation threshold. The landscape is rugged, and the trajectory is temporarily trapped in sharp, overfitting minima.
    
    \item \textbf{Descent 2 (Relaxation):} In the over-parameterized regime, the ``entropic selection'' mechanism of SGD dominates. The landscape-dependent noise ``anneals'' the system, driving it out of sharp minima and into flat minima. This \textit{implicit regularization} via non-equilibrium dynamics is the physical origin of the second descent.
\end{enumerate}

Mathematically, the escape rate from a sharp minimum with curvature $\lambda$ scales as:
\begin{equation}
    \Gamma_{\text{escape}} \sim \exp\left(-\frac{\Delta L}{D_{\text{eff}}}\right) \sim \exp\left(-\frac{\Delta L \cdot B}{\eta \lambda}\right)
\end{equation}
where $\Delta L$ is the barrier height. For sharp minima with large $\lambda$, the escape rate is high, explaining the instability of overfitting solutions.

\subsubsection{Predictions and Implications}

The NESP framework makes several testable predictions:

\begin{enumerate}[topsep=1pt,itemsep=1pt]
    \item \textbf{Batch size dependence:} Larger batch sizes reduce noise ($D \propto 1/B$), making sharp minima more stable and potentially suppressing the second descent.
    
    \item \textbf{Learning rate dependence:} Higher learning rates increase noise ($D \propto \eta$), accelerating escape from sharp minima and enhancing the second descent.
    
    \item \textbf{Epoch-wise double descent:} The non-monotonic behavior over training time is a natural consequence of the system first falling into sharp minima (fast convergence) then slowly escaping (relaxation).
\end{enumerate}

These predictions align with empirical observations in \cite{nakkiran_deep_2019} and provide a unified explanation for both model-wise and epoch-wise double descent phenomena.

\subsection{Experimental Verification: A Toy Model Proposal}

The theoretical framework developed in preceding sections---positing that Double Descent arises from a Non-Equilibrium phase transition between Jammed and Relaxed states, mediated by anisotropic SGD noise---demands empirical validation. However, the complexity of modern deep neural networks obscures the fundamental mechanisms we seek to isolate. Therefore, we propose a series of \textit{Gedankenexperimente} utilizing minimal, analytically tractable toy models that preserve the essential physics while eliminating confounding factors.

\subsubsection{The Teacher-Student Protocol}

We propose the \textbf{Teacher-Student framework} as our primary experimental testbed. This paradigm, extensively employed in the statistical physics of learning \cite{engel2001statistical,advani2013statistical}, provides the crucial advantage of a known ground truth, enabling precise measurement of generalization error without resort to heuristic train-test splits.

\paragraph{Setup.} Consider a \textit{Teacher network} $f^*(x) = \sigma(W^* x)$ with fixed, randomly drawn weights $W^* \in \mathbb{R}^{P^* \times d}$, where $\sigma(\cdot)$ denotes a nonlinearity (e.g., ReLU or $\tanh$). The Teacher generates labels:
\begin{equation}
    y = f^*(x) + \xi, \quad \xi \sim \mathcal{N}(0, \sigma_\xi^2)
\end{equation}
where $\xi$ represents irreducible label noise---a critical ingredient for inducing the interpolation peak.

The \textit{Student network} $\hat{f}(x; W) = \sigma(W x)$ with $W \in \mathbb{R}^{P \times d}$ is trained on $N$ samples $\{(x_i, y_i)\}_{i=1}^N$ drawn i.i.d. from the Teacher's distribution. By varying the Student's capacity $P$ while holding $N$ fixed, we sweep across the interpolation threshold $P/N \approx 1$.

\paragraph{The Critical Control Parameter.} The ratio $\alpha = P/N$ serves as our control parameter, analogous to the ``packing fraction'' in jamming physics. Our NESP framework predicts:
\begin{itemize}[topsep=1pt,itemsep=1pt]
    \item $\alpha \ll 1$: Under-parameterized regime. Classical bias dominated behavior.
    \item $\alpha \approx 1$: Critical/Jammed regime. Interpolation threshold.
    \item $\alpha \gg 1$: Over-parameterized regime. Entropic selection dominates.
\end{itemize}

\subsection{Physical Observables: Beyond Test Error}

To validate the NESP mechanism---rather than merely confirming the existence of Double Descent---we must measure \textit{physical observables} that directly probe the landscape geometry and noise structure. We propose three such observables:

\subsubsection{Observable I: The Hessian Trace (Sharpness)}

The trace of the Hessian, $\mathrm{Tr}(H) = \mathrm{Tr}(\nabla^2 L)$, provides a scalar summary of the local curvature at the solution $W^*_{\text{SGD}}$. This quantity is computationally accessible via Hutchinson's trace estimator:
\begin{equation}
    \mathrm{Tr}(H) \approx \frac{1}{K} \sum_{k=1}^K v_k^\top H v_k, \quad v_k \sim \mathcal{N}(0, I)
\end{equation}
where the Hessian-vector products $Hv$ can be computed efficiently via automatic differentiation.

\begin{hypothesis}[Sharpness Peak at Jamming]
The Hessian trace, measured at the end of training, exhibits a pronounced maximum at the interpolation threshold:
\begin{equation}
    \mathrm{Tr}(H)|_{\alpha \approx 1} \gg \mathrm{Tr}(H)|_{\alpha \ll 1} \quad \text{and} \quad \mathrm{Tr}(H)|_{\alpha \approx 1} \gg \mathrm{Tr}(H)|_{\alpha \gg 1}
\end{equation}
Furthermore, in the over-parameterized regime, we predict a \textbf{monotonic decay}:
\begin{equation}
    \frac{d}{d\alpha} \mathrm{Tr}(H) < 0 \quad \text{for} \quad \alpha > 1
\end{equation}
This decay directly evidences the system's migration from sharp (Jammed) to flat (Relaxed) minima as capacity increases.
\end{hypothesis}

\subsubsection{Observable II: Noise-Hessian Alignment}

The central claim of our NESP framework is that SGD noise is \textit{not} thermal (isotropic) but rather \textit{geometric} (aligned with the Hessian). To test this, we propose measuring the alignment between the noise covariance matrix $\Sigma(W)$ and the Hessian $H(W)$:
\begin{equation}
    \mathcal{A}(\Sigma, H) = \frac{\mathrm{Tr}(\Sigma H)}{\|\Sigma\|_F \|H\|_F}
\end{equation}
where $\|\cdot\|_F$ denotes the Frobenius norm. For thermal noise, $\Sigma = \sigma^2 I$ and thus $\mathcal{A} \propto \mathrm{Tr}(H)/\|H\|_F$, which is independent of the noise structure. For geometry-dependent noise, we expect strong positive correlation.

\begin{hypothesis}[Anisotropic Noise Alignment]
The alignment $\mathcal{A}(\Sigma, H)$ is significantly elevated compared to the null hypothesis of isotropic noise:
\begin{equation}
    \mathcal{A}(\Sigma, H) \gg \mathcal{A}(\sigma^2 I, H) \quad \text{across all } \alpha
\end{equation}
Moreover, this alignment should \textbf{strengthen} near the interpolation threshold, where the curvature structure is most pronounced. This would provide direct evidence that SGD performs geometric, not thermal, exploration of the loss landscape.
\end{hypothesis}

\subsubsection{Observable III: Eigenspectrum Evolution}

The full eigenspectrum of the Hessian, $\{\lambda_i\}_{i=1}^P$, encodes richer information than the trace alone. We propose tracking the spectral density $\rho(\lambda) = (1/P)\sum_i \delta(\lambda - \lambda_i)$ across the $\alpha$-sweep.

\begin{hypothesis}[Spectral Signature of Phase Transition]
At the Jamming transition ($\alpha \approx 1$), the spectral density exhibits:
\begin{enumerate}[topsep=1pt,itemsep=1pt]
    \item A \textbf{hard edge} at $\lambda = 0$, corresponding to the vanishing of the null space as constraints saturate degrees of freedom.
    \item \textbf{Outlier eigenvalues} at large $\lambda$, corresponding to sharp directions associated with fitting noise.
\end{enumerate}
In the Relaxed phase ($\alpha \gg 1$), the spectrum should exhibit:
\begin{enumerate}[topsep=1pt,itemsep=1pt]
    \item An extensive \textbf{null space} (bulk of eigenvalues at $\lambda \approx 0$), corresponding to flat directions in the solution manifold.
    \item Suppression of outliers, indicating the system has escaped sharp directions.
\end{enumerate}
\end{hypothesis}

\subsection{Experimental Protocol}

To empirically validate the Entropic Selection mechanism, we propose the following experimental protocol:

\begin{algorithm}[htb]
\renewcommand{\algorithmcfname}{Experimental Protocol}
\caption{NESP Validation via Teacher-Student Model}
\KwData{Teacher weights $W^*$, noise level $\sigma_\xi$, sample sizes $N$, capacity range $P \in [P_{\min}, P_{\max}]$.}
\KwResult{Measurements of $R_{\text{test}}(\alpha)$, $\mathrm{Tr}(H)(\alpha)$, $\mathcal{A}(\Sigma, H)(\alpha)$, $\rho(\lambda; \alpha)$.}
\Begin{
    \For{$P \in \{P_{\min}, \ldots, P_{\max}\}$}{
        Initialize Student network with random weights $W_0$\;
        Train via SGD: $W_{t+1} = W_t - \eta \nabla L_{\mathcal{B}}(W_t)$\;
        \tcp*[l]{At convergence:}
        Compute $R_{\text{test}}(W_T) = \mathbb{E}_{x}[(f^*(x) - \hat{f}(x; W_T))^2]$\;
        Compute $\mathrm{Tr}(H)$ via Hutchinson estimator\;
        Estimate $\Sigma(W_T)$ from mini-batch gradient variance\;
        Compute alignment $\mathcal{A}(\Sigma, H)$\;
        (Optional) Compute top-$k$ Hessian eigenvalues via Lanczos iteration\;
    }
}
\Return{Observables as functions of $\alpha = P/N$.}
\end{algorithm}

\paragraph{Predicted Signatures.} If our NESP framework is correct, the experimental results should exhibit the following correlated signatures:
\begin{enumerate}[topsep=1pt,itemsep=1pt]
    \item $R_{\text{test}}(\alpha)$ displays the characteristic Double Descent curve.
    \item $\mathrm{Tr}(H)(\alpha)$ peaks at $\alpha \approx 1$ and decays for $\alpha > 1$, \textit{anti-correlated} with the second descent.
    \item $\mathcal{A}(\Sigma, H)$ remains elevated throughout, peaking near $\alpha \approx 1$.
    \item The spectral density transitions from hard-edge (Jammed) to bulk-zero (Relaxed) as $\alpha$ crosses unity.
\end{enumerate}

The simultaneous observation of these four signatures would constitute strong evidence for the NESP interpretation of Double Descent as a dynamical phase transition driven by geometry-dependent noise.

\subsubsection{Alternative Minimal Model: Two-Layer Linear Networks}

For maximal analytical tractability, we also propose the \textbf{two-layer linear network}:
\begin{equation}
    \hat{f}(x; U, V) = U V^\top x, \quad U \in \mathbb{R}^{d_{\text{out}} \times k}, \; V \in \mathbb{R}^{d_{\text{in}} \times k}
\end{equation}
Despite being linear in the input-output map (equivalent to $W = UV^\top$), this model exhibits \textit{nonlinear dynamics} under gradient descent due to the factorization. Recent work \cite{advani2017highdimensionaldynamicsgeneralizationerror,saxe2014exact} has shown that this model captures essential features of deep learning dynamics, including staged learning and critical slowdowns.

The advantage of this model is that the Hessian and noise covariance can be computed \textit{analytically} in certain regimes, enabling precise theoretical predictions against which experimental measurements can be compared. The NESP predictions for this model are:
\begin{itemize}[topsep=1pt,itemsep=1pt]
    \item The ``effective temperature'' $T_{\text{eff}} = \eta \cdot \mathrm{Tr}(\Sigma)/P$ should exhibit non-monotonic behavior, peaking at the interpolation threshold.
    \item The condition number $\kappa(H) = \lambda_{\max}/\lambda_{\min}$ should diverge as $\alpha \to 1$, signaling the Jamming transition.
\end{itemize}

\subsection{Extension}

The idea itself offers many extensions for deployments.

\subsubsection{Synthesis: From Equilibrium to Non-Equilibrium}

In this work, we have undertaken a critical re-examination of the theoretical foundations underlying the Double Descent phenomenon. Our central thesis is that the apparent ``anomaly'' of Double Descent---its violation of the classical bias-variance tradeoff---is not a failure of statistical physics \textit{per se}, but rather a failure of \textbf{equilibrium} statistical physics inappropriately applied to an inherently \textbf{non-equilibrium} dynamical process.

We have argued that the dominant theoretical approaches---Random Matrix Theory, Replica Methods, and classical bias-variance decomposition---share a common, often implicit, assumption: that the learning algorithm converges to a stationary state characterized by a Boltzmann-like distribution over weight space. This equilibrium assumption, while mathematically convenient, fundamentally mischaracterizes the nature of Stochastic Gradient Descent.

Our reformulation, grounded in Non-Equilibrium Statistical Physics (NESP), makes three key departures from the equilibrium paradigm:

\begin{enumerate}[topsep=1pt,itemsep=1pt]
    \item \textbf{From static distributions to dynamical evolution:} We model the probability density $\rho(W,t)$ via the Fokker-Planck equation, treating the ``solution'' as a transient state rather than an equilibrium configuration.
    
    \item \textbf{From isotropic to anisotropic noise:} We recognize that SGD noise is not thermal but geometric---its covariance $\Sigma(W)$ is state-dependent and aligned with the loss landscape curvature $H(W)$.
    
    \item \textbf{From energy minimization to entropic selection:} We posit that generalization arises not from finding the global minimum, but from the \textit{dynamical stability} of flat minima under geometry-dependent noise. Sharp minima, despite potentially lower loss, are rendered unstable by their own curvature.
\end{enumerate}

\subsubsection{Bridging the Gap: The Play and The Stage}

A persistent tension in theoretical machine learning has been the disconnect between two perspectives:
\begin{itemize}[topsep=1pt,itemsep=1pt]
    \item \textbf{The Stage (Landscape):} Static analysis of the loss surface---its minima, saddle points, and global structure.
    \item \textbf{The Play (Dynamics):} The trajectory of the optimizer through weight space---its convergence, oscillations, and implicit biases.
\end{itemize}

Equilibrium theories privilege the Stage: they characterize the set of possible solutions and assume the optimizer will find them. But they cannot explain \textit{which} solution is selected, nor \textit{how} the selection depends on hyperparameters like learning rate and batch size.

The NESP framework bridges this gap by coupling the Stage to the Play. The landscape geometry (Hessian) determines the noise structure (diffusion tensor), which in turn shapes the dynamical trajectory. Double Descent emerges naturally from this coupling: the interpolation peak corresponds to a \textit{Jamming transition} where landscape constraints saturate available degrees of freedom, and the second descent corresponds to \textit{entropic relaxation} as the system escapes sharp minima driven by its own noise.

This perspective offers a resolution to the longstanding puzzle: \textit{Why do over-parameterized models generalize?} Our answer: not because they find better minima (they often do not), but because SGD dynamics are \textit{geometrically biased} toward flat regions of the solution manifold. Generalization is an emergent property of non-equilibrium dynamics, not a static property of the loss landscape.

Moreover, if the NESP framework, is formally validated, it opens several exciting avenues for future research, as followed. 

\subsubsection{Grokking as Extreme Relaxation}

The phenomenon of ``Grokking''---wherein models exhibit delayed generalization long after achieving zero training error \cite{Power2022Grokking}---presents a striking challenge to equilibrium theories. From the NESP perspective, we speculate that Grokking represents an \textit{extreme} instance of the Jamming-Relaxation transition.

Specifically, certain data structures (e.g., modular arithmetic) may induce loss landscapes where the Jammed phase is \textit{metastable} over extended timescales. The system achieves interpolation (zero training loss) but remains trapped in sharp, non-generalizing minima. Generalization occurs only after sufficient time for the noise-driven annealing process to escape these metastable states.

This interpretation predicts that Grokking should be:
\begin{itemize}[topsep=1pt,itemsep=1pt]
    \item \textbf{Accelerated} by increasing learning rate or decreasing batch size (higher effective noise).
    \item \textbf{Suppressed} by explicit regularization that pre-selects flat minima.
    \item \textbf{Accompanied} by a monotonic decrease in Hessian trace during the ``grokking'' phase.
\end{itemize}

\subsubsection{Transformers and Large Language Models}

The training dynamics of Large Language Models (LLMs) represent perhaps the most consequential---and least understood---application of deep learning. The NESP framework suggests several speculative connections:

\paragraph{Temperature in Sampling vs. Training.} LLMs employ a ``temperature'' parameter $T$ during inference to control the sharpness of the output distribution. Intriguingly, our framework identifies an ``effective temperature'' $T_{\text{eff}} = \eta \cdot \mathrm{Tr}(\Sigma)/(2B)$ that governs training dynamics. We pose the question: \textit{Is there a principled relationship between these two temperatures?}

One speculative hypothesis: models trained with higher $T_{\text{eff}}$ (larger learning rate, smaller batch) may naturally produce sharper internal representations, requiring lower inference temperature for coherent outputs. This would establish a non-trivial link between training dynamics and inference behavior.

\paragraph{Scaling Laws and Phase Transitions.} Empirical ``scaling laws'' \cite{kaplan2020scalinglawsneurallanguage} describe how LLM performance improves with model size, data, and compute. From the NESP perspective, these smooth power-law scalings may obscure underlying \textit{phase transitions} in the training dynamics. As model capacity crosses critical thresholds relative to data complexity, the system may undergo Jamming-Relaxation transitions analogous to Double Descent.

Understanding these transitions could inform more efficient training schedules: rather than smooth scaling of learning rate, one might employ \textit{punctuated} schedules that exploit phase boundaries---high noise to escape Jammed phases, low noise to stabilize in Relaxed phases.

\subsubsection{Toward a Unified Theory of Implicit Regularization}

The NESP framework provides a candidate mechanism for ``implicit regularization''---the empirical observation that SGD finds generalizing solutions without explicit regularization terms. Our entropic selection mechanism offers a \textit{dynamical} explanation: SGD implicitly regularizes by destabilizing sharp minima.

A unified theory would connect this mechanism to other forms of implicit bias:
\begin{itemize}[topsep=1pt,itemsep=1pt]
    \item \textbf{Minimum-norm solutions} in linear regression.
    \item \textbf{Max-margin classifiers} in linearly separable problems.
    \item \textbf{Low-rank bias} in matrix factorization.
\end{itemize}

We conjecture that all these biases arise from the same underlying principle: the geometry-dependent noise of SGD preferentially stabilizes solutions with specific structural properties (low norm, large margin, low rank) because these properties correlate with flat loss landscape geometry.

